\rtmtight
\begin{tblr}{width=\textwidth,colspec={|Q[c,m,2.1cm]|X[l,m]|},hlines,vlines,hspan=minimal,rowsep=1.5pt,colsep=3pt,row{1}={bg=headgray,font=\bfseries}}
\SetCell{c,m}\hfil\logo\hfil & \textbf{POLITEKNIK NEGERI MALANG}\par\textbf{JURUSAN TEKNOLOGI INFORMASI}\par\textbf{PROGRAM STUDI: D4 TEKNIK INFORMATIKA} \\
\SetCell[c=2]{l} \textbf{RENCANA TUGAS MAHASISWA (RTM) DAN RUBRIK PENILAIAN} & \\
Nama Mata Kuliah & Proyek Teknologi Terintegrasi \\
Kode Mata Kuliah & RTI256104 \quad \textbf{sks / Jam perminggu:} 6 SKS/12 jam \quad \textbf{Semester:} 6 \\
Dosen Pengampu & Tim Pengajar Program Studi D4 TI \\
\end{tblr}
\assessmentblock{Dokumen Perencanaan Proyek TI Terintegrasi}{Case Method dan penyusunan dokumen proyek}{SCPMK0801-04801}{Mahasiswa menyusun dokumen perencanaan proyek TI terintegrasi yang mencakup WBS, jadwal, estimasi biaya, alokasi sumber daya, dan risk register secara komprehensif.}{Menganalisis kebutuhan proyek, menyusun WBS dan jadwal Gantt, mengestimasi anggaran, mengidentifikasi risiko, dan mempresentasikan dokumen kepada tim penilai.}{Dokumen perencanaan proyek (WBS, Gantt chart, anggaran, risk register) dan bahan presentasi.}{Kelengkapan WBS, jadwal, dan hierarki pekerjaan & 6 & WBS lengkap sampai level aktivitas, jadwal Gantt realistis dengan milestone, dependensi, dan jalur kritis teridentifikasi. & WBS dan jadwal tersedia tetapi beberapa aktivitas, dependensi, atau milestone belum lengkap. & WBS tidak mencapai level aktivitas atau jadwal tidak dapat digunakan sebagai acuan proyek. \\ \\
Kualitas estimasi anggaran dan risk register & 5 & Anggaran dirinci per WBS item dengan justifikasi, risiko diidentifikasi dengan probabilitas, dampak, dan rencana mitigasi. & Anggaran dan risiko tersedia tetapi justifikasi atau rencana mitigasi masih terbatas. & Anggaran tidak dirinci atau risiko tidak memiliki rencana mitigasi yang operasional. \\ \\
Presentasi dokumen dan kesiapan sprint & 4 & Dokumen dipresentasikan secara sistematis, rencana sprint awal jelas, dan tim siap memulai implementasi. & Presentasi cukup jelas tetapi rencana sprint awal atau kesiapan tim masih kurang rinci. & Presentasi tidak dapat menjelaskan rencana proyek atau tim belum siap memulai sprint. \\}

\assessmentblock{Sprint Implementasi Multi-Platform (Sprint 1-4)}{Project Based Learning}{SCPMK0802-04802, SCPMK0607-04805}{Mahasiswa mengimplementasikan aplikasi multi-platform secara bertahap melalui empat sprint agile dengan dokumentasi lengkap, integrasi komponen, dan kolaborasi tim yang terukur.}{Melaksanakan sprint planning, mengimplementasikan fitur sesuai backlog, melakukan code review dan integrasi, mendokumentasikan sprint log, dan mempresentasikan progres.}{Kode sumber terintegrasi di repositori, sprint log, burndown chart, dokumentasi API, dan demo produk per sprint.}{Kualitas implementasi fitur dan kepatuhan sprint & 16 & Fitur terdeliver sesuai sprint goal, kode berjalan tanpa error kritis, acceptance criteria terpenuhi, dan backlog terkelola. & Sebagian besar fitur terdeliver tetapi beberapa acceptance criteria atau kualitas kode belum optimal. & Fitur tidak terdeliver sesuai sprint goal atau kode mengandung error kritis yang tidak diselesaikan. \\ \\
Integrasi front-end, back-end, dan database & 14 & Komponen terintegrasi dengan API yang terdokumentasi, database schema konsisten, dan alur data berjalan end-to-end. & Integrasi sebagian besar berjalan tetapi beberapa endpoint, schema, atau alur data masih bermasalah. & Integrasi tidak berjalan atau komponen tidak dapat berkomunikasi secara konsisten. \\ \\
Dokumentasi sprint dan kolaborasi tim & 10 & Sprint log, burndown chart, commit history, dan pembagian tugas terdokumentasi lengkap serta menunjukkan kolaborasi aktif. & Dokumentasi sprint tersedia tetapi beberapa komponen atau bukti kolaborasi masih kurang. & Dokumentasi sprint tidak lengkap atau tidak ada bukti kolaborasi tim yang bermakna. \\}

\assessmentblock{UTS Milestone Review Sprint 1-2}{UTS berbasis presentasi dan demo}{SCPMK0803-04803}{Mahasiswa mempresentasikan progres proyek sprint 1-2, mendemonstrasikan produk yang sudah berjalan, menyajikan laporan kemajuan, dan mengelola backlog untuk sprint berikutnya.}{Demo produk sprint 1-2, presentasi burndown chart dan velocity, penyajian backlog terupdate, dan menjawab pertanyaan reviewer.}{Demo produk berjalan, laporan kemajuan proyek, backlog terupdate, dan bahan presentasi milestone.}{Kualitas demo produk sprint 1-2 & 7 & Demo menampilkan fitur yang telah selesai, berjalan tanpa error, dan memenuhi acceptance criteria sprint 1-2. & Demo berjalan tetapi beberapa fitur tidak stabil atau acceptance criteria belum sepenuhnya terpenuhi. & Demo tidak dapat berjalan atau fitur yang ditunjukkan tidak sesuai sprint goal. \\ \\
Pengelolaan backlog dan monitoring progres & 5 & Backlog terkelola dengan prioritas yang jelas, velocity dihitung, burndown chart konsisten dengan progres aktual. & Backlog tersedia dan diperbarui tetapi prioritas, velocity, atau burndown chart belum optimal. & Backlog tidak terkelola atau tidak ada bukti monitoring progres yang terstruktur. \\ \\
Laporan kemajuan dan komunikasi proyek & 3 & Laporan kemajuan mencakup status, risiko aktual, tindakan korektif, dan rencana sprint berikutnya secara jelas. & Laporan kemajuan tersedia tetapi analisis risiko atau rencana tindakan korektif masih kurang. & Laporan tidak mencerminkan kondisi aktual proyek atau tidak ada rencana sprint berikutnya. \\}

\assessmentblock{UAS Defense Proyek Akhir}{UAS berbasis defense dan demo}{SCPMK0804-04804, SCPMK0803-04803}{Mahasiswa mempresentasikan produk akhir proyek TI terintegrasi yang siap deploy, melakukan defense teknis dan manajerial secara komprehensif, serta menyerahkan laporan akhir proyek.}{Demo produk deployable, presentasi laporan akhir, menjawab pertanyaan teknis dan manajerial dari tim penilai secara individual maupun tim.}{Produk akhir deployable, laporan akhir proyek lengkap, repository kode, dan bahan presentasi defense.}{Kualitas produk akhir dan deployability & 12 & Produk berjalan di environment produksi/staging, semua fitur utama berfungsi, dan dapat digunakan oleh pengguna nyata. & Produk sebagian besar berfungsi tetapi beberapa fitur tidak stabil atau deployment belum lengkap. & Produk tidak dapat dijalankan di luar development environment atau banyak fitur tidak berfungsi. \\ \\
Kelengkapan laporan akhir proyek & 10 & Laporan mencakup arsitektur, implementasi, pengujian, hasil UAT, deployment guide, dan retrospektif proyek secara lengkap. & Laporan tersedia tetapi beberapa komponen seperti hasil UAT, deployment guide, atau retrospektif masih kurang. & Laporan tidak lengkap atau tidak mencerminkan produk yang dibuat secara akurat. \\ \\
Defense teknis dan kemampuan argumentasi & 8 & Mahasiswa mempertahankan keputusan teknis dan manajerial dengan evidence, menjelaskan trade-off, dan merespons pertanyaan secara kompeten. & Defense cukup baik tetapi beberapa keputusan teknis atau manajerial belum dapat dipertahankan dengan evidence. & Mahasiswa tidak dapat menjelaskan keputusan desain atau memberikan respons yang relevan terhadap pertanyaan. \\}

\begin{tblr}{width=\textwidth,colspec={|Q[l,m,3.2cm]|X[l,m]|},hlines,vlines,hspan=minimal,row{1-3}={bg=headgray},rowsep=1.5pt,colsep=3pt}
Jadwal Pelaksanaan: & Minggu 1--17 sesuai RPS. Setiap evaluasi memiliki RTM dan rubrik multi-kriteria. \\
Lain-Lain Yang Diperlukan: & LMS, repositori proyek, perangkat lunak sesuai mata kuliah, dan perangkat presentasi. \\
Pustaka: & Mengacu pustaka utama dan pendukung pada bagian RPS. \\
\end{tblr}
